\scanpage{37}\mspage{36}

et en fait, le contraire est vrai dans les \inserted{\unsure{2|c$^{\mathrm t}$}} exemples 1. A) et B).

\uline{Définition}\quad On dit que $\unsure{\mathcal F|\mathscr F}$ est un \uline{atelier ensembliste} si pour $F,G\in\mathcal F$, on a
\[
\begin{aligned}
&\text{a) }\mathcal F\longrightarrow\mathrm{Fig}(L)\quad\text{injectif}\\
&\text{b) }\mathrm{mulomb}\,F\leq\mathrm{mulomb}\,G\Longrightarrow F\leq G\\
&\text{c) }\mathrm{mulomb}\,F\ll\mathrm{mulomb}\,G\Longrightarrow F\ll G\\
&\deleted{\unsure{\mathrm{d)}}}\ \unsure{\mathrm{d)}}\,\mathrm{mulomb}\,F\mathrel{\tikz[baseline=-.55ex,x=.85em,y=.70em]{\draw[line cap=round,line join=round] (0,1) .. controls (.22,.77) and (.86,.67) .. (.93,.61) .. controls (.82,.32) and (.28,.16) .. (.10,0); \draw[line cap=round,line join=round] (.98,1) .. controls (.77,.82) and (.43,.50) .. (.26,.37) .. controls (.43,.32) and (.72,.18) .. (.97,.04);}}\mathrm{mulomb}\,G\Longrightarrow F\mathrel{\tikz[baseline=-.55ex,x=.85em,y=.70em]{\draw[line cap=round,line join=round] (0,1) .. controls (.22,.77) and (.86,.67) .. (.93,.61) .. controls (.82,.32) and (.28,.16) .. (.10,0); \draw[line cap=round,line join=round] (.98,1) .. controls (.77,.82) and (.43,.50) .. (.26,.37) .. controls (.43,.32) and (.72,.18) .. (.97,.04);}}G\deleted{\unsure{\mathcal F}}\\
&\deleted{\mathrm{mulomb}(F)\ \text{élémentaire}\ \Longleftrightarrow}
\end{aligned}
\]
% slanted
\margin{De \unsure{toutes façons|toute façon}, un axiome ultérieur, imposé pour \uline{\unsure{tout}} atelier, \unsure{indiquerait|indiquera} que \deleted{\unsure{pour}} $\mathrm{mulomb}(F)$ élémentaire $\Longleftrightarrow F$ élémentaire}

(NB peut-être qu'on sera amené encore à renforcer ces conditions), de telle sorte que
\[
\mathcal F\longrightarrow\mathrm{Fig}(L)
\]
identifie $\mathcal F$, avec les structures $\leq,\ll,\mathrel{\tikz[baseline=-.55ex,x=.85em,y=.70em]{\draw[line cap=round,line join=round] (0,1) .. controls (.22,.77) and (.86,.67) .. (.93,.61) .. controls (.82,.32) and (.28,.16) .. (.10,0); \draw[line cap=round,line join=round] (.98,1) .. controls (.77,.82) and (.43,.50) .. (.26,.37) .. controls (.43,.32) and (.72,.18) .. (.97,.04);}}$, à un \deleted{sous-ensemble} ``sous-atelier'' de $\mathrm{Fig}(L)$.

Je voudrais préciser ce point, \deleted{en termes de} \inserted{D'abord}

\uline{Définition}\quad Soient $(\mathcal F,\leq,\ll)$ et $(\mathcal F',\leq,\ll)$ deux ateliers, $f=\unsure{f_{\mathrm{fig}}^{\prime}|f_{\mathrm{fig}}}\colon\mathcal F\longrightarrow\mathcal F'$ une application. On dit que $f$ est un \uline{morphisme d'ateliers} si

a) $f$ croissante pour $\leq$ et $\ll$

b) \deleted{$f$ transforme \unsure{sous}} $f$ commute aux $\mathrm{Sup}$ pour $\leq$ \inserted{(\unsure{d'où} $\varnothing_{\mathcal F}\mapsto\varnothing_{\mathcal F'}$)}, et aux $\mathrm{Inf}$ \inserted{\deleted{\unsure{major|non}} des familles \unsure{majorées|n.v. paires}} pour $\leq$

\deleted{\unsure{c)} $f$ transforme figures élémentaires} \deleted{\unsure{Cela|(cela} revient au $=$ que de se donner} 

% margin note written sideways
\margin{En fait, si $F'=f(F)$ alors $\mathcal F_{\leq F}\simeq{\unsure{\mathcal P|\mathfrak P}}_{\mathrm{fermés}}(\widetilde F)$ $\mathcal F'_{\leq F'}\simeq\mathcal P_{\unsure{\mathrm{fermé}|\mathrm{fermés}}}(\widetilde F')$, et l'application $(*)$\tikz[remember picture,overlay]\coordinate (ms37UpperNote); \unsure{semble} être \unsure{induite|induit} \unsure{tjs|tjrs} par $\widetilde F'\unsure{\rule[.5ex]{1.35em}{.6pt}|\longrightarrow}\widetilde F$ \uline{croissante} \inserted{i.e. continue} et $=$ telle que \unsure{tt|le} \unsure{fermé|suppt} de $\widetilde F'$ soit image inv. de un \unsure{fermé} de $\widetilde F$\,\ldots}
% slanted
\margin{\tikz[remember picture,overlay]\coordinate (ms37PairOne);\uline{NB} a) Dans tous les exemples que j'ai vus, $\deleted{\unsure{A|\forall}}\forall\,F\in\mathcal F$, l'application induite par $f$ $(*)\tikz[remember picture,overlay]\coordinate (ms37LowerNote);\ \mathcal F_{\leq F}\longrightarrow\mathcal F'_{\leq f(F)}$ est \uline{surjective}. \tikz[remember picture,overlay]\coordinate (ms37PairTwo);b) Dans la plupart des cas (``cas ensembliste'', \unsure{notamment|notam$^{\mathrm t}$}) $f$ est croissante pour $\unsure{\preceq|\leq}$ également\,\ldots}
\figure{\tikz[remember picture,overlay]{
  \draw[->] ([xshift=-2pt,yshift=1pt]ms37LowerNote) .. controls +(13pt,14pt) and +(4pt,-12pt) .. ([xshift=-2pt,yshift=-1pt]ms37UpperNote);
  \draw ([xshift=-3pt,yshift=5pt]ms37PairOne) -- ([xshift=-3pt,yshift=-6pt]ms37PairOne);
  \draw ([xshift=-4.5pt,yshift=5pt]ms37PairOne) -- ([xshift=-4.5pt,yshift=-6pt]ms37PairOne);
  \draw ([xshift=-3pt,yshift=5pt]ms37PairTwo) -- ([xshift=-3pt,yshift=-7pt]ms37PairTwo);
  \draw ([xshift=-4.5pt,yshift=5pt]ms37PairTwo) -- ([xshift=-4.5pt,yshift=-7pt]ms37PairTwo);
}}

$\deleted{\unsure{f_{\mathrm{str}}}\colon M\longrightarrow M'}$
On dit \deleted{telle que} que $f$ est ``compatible aux strates'' si $f$ satisfait \deleted{\unsure{a)} $f$ croissante pour} en plus

c) $f$ transforme multistrates \inserted{(i.e. fig. él.)} en multistrates \inserted{(fig. él.)}

Alors la donnée d'une telle $f$ équivaut à celle de
\[
\unsure{f_{\mathrm{strat}}|f_{\mathrm{str.t}}}\colon M\unsure{\rule[.5ex]{1.65em}{.65pt}|\longrightarrow}M'
\]
